\documentclass[11pt]{article}
\usepackage[a4paper,margin=26mm]{geometry}
\usepackage[T1]{fontenc}
\usepackage{lmodern,amsmath,amssymb,amsthm,mathtools,microtype,booktabs,array,enumitem}
\usepackage[hidelinks]{hyperref}
\setlength{\emergencystretch}{3em}
\setlist{nosep,leftmargin=*}
\numberwithin{equation}{section}
\newtheorem{theorem}{Theorem}[section]
\newtheorem{proposition}[theorem]{Proposition}
\newtheorem{lemma}[theorem]{Lemma}
\newtheorem{corollary}[theorem]{Corollary}
\theoremstyle{definition}
\newtheorem{assumption}[theorem]{Assumption}
\theoremstyle{remark}
\newtheorem{remark}[theorem]{Remark}
\newcommand{\R}{\mathbb R}
\newcommand{\E}{\mathbb E}
\newcommand{\PP}{\mathbb P}
\newcommand{\Law}{\operatorname{Law}}
\newcommand{\Cov}{\operatorname{Cov}}
\newcommand{\supp}{\operatorname{supp}}
\newcommand{\one}{\mathbf 1}
\newcommand{\Ga}{\Gamma}
\newcommand{\ee}{\mathrm e}
\newcommand{\ii}{\mathrm i}
\newcommand{\dd}{\,\mathrm d}
\newcommand{\BS}{Backhausz--Szegedy}
\title{\textbf{Normalized Gaussian waves for bulk eigenvectors}\\[3pt]
\large A local extension of Backhausz--Szegedy, a reduced bulk-universality input,\\
and a proof structure for later Lean formalization}
\author{Research note on the two supplied manuscripts}
\date{November 4, 2026}
\begin{document}
\maketitle
\begin{abstract}
The useful extension is a uniform, conditional-on-the-graph Gaussian-wave
limit with variance one for fixed orthonormal tuples at nearby energies.
We prove this on every fixed tree neighborhood, assuming only a uniform
scalar fourth-moment bound in probability. The published input is
\BS's Theorem 2.2 and Remark 1; their scalar Theorem 2.1 alone is insufficient
for the local joint conclusion. We then identify a smaller quantitative
input sufficient for the supplied bulk-universality argument: covariance
control for quadratic marks of two eigenvectors on one oriented edge.
This input implies the required scalar moment bound and the weighted
spectral-tail estimate. We distinguish the results proved here from the
additional quantitative theorem asserted and argued in the supplied draft.
\end{abstract}

\section{What the source article proves and how its proof works}
The source is \BS, \emph{On the almost eigenvectors of random regular
graphs}, Ann. Probab.\ 47 (2019), 1677--1725 \cite{BS}. The second source
is the supplied November 3 bulk-universality TeX draft \cite{Bulk}.
All references to numbered results of \BS\ below use the published PDF.

\begin{center}
\small
\begin{tabular}{>{\raggedright\arraybackslash}p{.25\linewidth}>{\raggedright\arraybackslash}p{.48\linewidth}>{\raggedright\arraybackslash}p{.16\linewidth}}
\toprule
Source result & Content & PDF / journal page\\
\midrule
Theorem 2.1 & Uniform weak Gaussianity of the empirical entries of every
zero-sum almost eigenvector; the limiting variance may be below one.
& 4 / 1680\\[3pt]
Corollary 2.2 & Joint Gaussian entry laws for fixed tuples at nearby
approximate eigenvalues. & 5 / 1681\\[3pt]
Theorem 2.2, Remark 1 & Random-covering neighborhood laws; the remark
extends the conclusion to fixed tuples at nearby energies.
& 6 / 1682\\
\bottomrule
\end{tabular}
\end{center}

Their conclusion allows a scalar wave $\sigma\Psi_\lambda$ with
$0\leq\sigma\leq1$. Although the finite vectors have normalized squared
length, weak convergence can lose second moments. Thus ``exact
eigenvector'' and ``orthonormal'' do not, on their own, justify replacing
$\sigma$ by one in this argument.

The proof of Theorem 2.1 is obtained from the stronger local theorem.
Its main dependency chain is as follows.
\begin{enumerate}
\item Proposition 5.1 transfers the classification of typical tree
processes, Theorem 5.2, to uniform finite-graph statements. Compactness,
concentration of local-global colored laws, and a contradiction argument
produce common graph events that work for all admissible vectors.
\item Lemma 5.3 uses spectral bounds and truncation to put a nonconstant
typical eigenvector process in the tree spectral interval.
\item Theorem 6.1 gives a finite-alphabet entropy inequality by counting
colored regular graphs. Proposition 7.1 and Appendix C permit smoothing
with an independent Gaussian wave while preserving typicality.
\item Theorems 8.1 and 9.1 compare enlarged-neighborhood subspace
entropies with the Gaussian benchmark. Monotonicity, summarized in
Theorem 10.3, gives the lower benchmark at one star; equality also gives
the 2-Markov property.
\item Theorem 11.2 gives the upper benchmark, with equality only at the
Gaussian star. Its proof uses Gaussian convolution, the de Bruijn
identity, Fisher information, and a weighted frame identity.
\item The two bounds force equality. The Gaussian star and the
2-Markov property determine the whole Gaussian tree process. This
completes Theorem 5.2 and hence Theorem 2.2.
\end{enumerate}
This compactness-based proof supplies no polynomial rate. A quantitative
eigenvector-moment theorem needs additional estimates; see
Section~\ref{sec:quantitative}.

\section{Sampling, normalization, and exact local covariance}
Fix $d\geq3$ and a compact interval
$K\Subset(-2\sqrt{d-1},2\sqrt{d-1})$. Let $G$ be a simple $d$-regular
graph on $N$ vertices, with adjacency matrix $A$. Choose real
orthonormal eigenvectors
\[
 Au_a=\lambda_a u_a,\qquad
 \langle u_a,u_b\rangle=\one_{a=b},\qquad \lambda_a\in K.
\]
Since $\lambda_a\ne d$ and $A\one=d\one$, each $u_a$ has zero entry sum.
No connectedness or simplicity of the eigenvalues is required.

Let $T_d$ be the infinite regular tree with root $o$. Sample a covering
map $\phi:T_d\to G$ by choosing $\phi(o)$ uniformly, mapping its
neighbors by a uniform bijection, and continuing by independent uniform
local bijections away from the root. For finite $F\subset V(T_d)$ set
\begin{equation}\label{eq:X}
 X^F=(X_{a,x})_{a\leq k,x\in F},\qquad
 X_{a,x}=\sqrt N\,u_a(\phi(x)).
\end{equation}
Write $\E_\phi$ for this sampling expectation with the graph and vectors
fixed. Every $\phi(x)$ is uniform on the graph. The lifted law is
invariant under tree automorphisms. Distinct tree vertices can have the
same image; the covering law is defined even if $G$ contains short cycles.

Define the spherical covariance polynomials recursively by
\begin{equation}\label{eq:recurrence}
 \varphi_\lambda(0)=1,\qquad
 \varphi_\lambda(1)=\lambda/d,\qquad
 (d-1)\varphi_\lambda(r+1)=\lambda\varphi_\lambda(r)-\varphi_\lambda(r-1)
 \quad(r\geq1).
\end{equation}
For the variance-one Gaussian wave $\Psi_\lambda$,
$\E\Psi_\lambda(x)\Psi_\lambda(y)=\varphi_\lambda(\operatorname{dist}(x,y))$.
Denote its covariance on $F$ by $\Ga_F(\lambda)$ and let
\begin{equation}\label{eq:target}
 \gamma^F_{\boldsymbol\lambda}
 =\bigotimes_{a=1}^k N(0,\Ga_F(\lambda_a)).
\end{equation}
These are measures on the physical coordinates $\R^{k|F|}$; their
covariance may be singular because of eigenvector relations.

\begin{lemma}[Exact covering covariance]\label{lem:cov}
On every fixed graph as above,
\begin{equation}\label{eq:cov}
 \E_\phi X_{a,x}=0,\qquad
 \E_\phi X_{a,x}X_{b,y}
 =\one_{a=b}\varphi_{\lambda_a}(\operatorname{dist}(x,y)).
\end{equation}
\end{lemma}
\begin{proof}
Uniformity of $\phi(x)$ gives the mean and the distance-zero covariance
$\langle u_a,u_b\rangle$. By invariance the covariance depends only on
$r=\operatorname{dist}(x,y)$; write it $c_{ab}(r)$. The eigenvector
relation at $y$ gives
\[
 d\,c_{ab}(1)=\lambda_b\,c_{ab}(0),\qquad
 c_{ab}(r-1)+(d-1)c_{ab}(r+1)=\lambda_b c_{ab}(r),\quad r\geq1.
\]
There is one neighbor of $y$ nearer $x$ and $d-1$ farther away.
The initial values and recurrence uniquely give \eqref{eq:cov}.
\end{proof}

An equivalent finite-matrix proof uses nonbacktracking-walk matrices
\[
 Q_0(A)=I,\quad Q_1(A)=A,\quad Q_2(A)=A^2-dI,\quad
 Q_r(A)=AQ_{r-1}(A)-(d-1)Q_{r-2}(A)\quad(r\geq3).
\]
For $r\geq1$ there are $d(d-1)^{r-1}$ such walks from each vertex.
Their endpoint covariance is
$\langle u_a,Q_r(A)u_b\rangle/[d(d-1)^{r-1}]$.
The exceptional formula for $Q_2$ must be retained; applying the later
recurrence already at $r=2$ would give an incorrect diagonal term.

For the star $F=\{o,o_1,\ldots,o_d\}$, the one-channel covariance is
\begin{equation}\label{eq:star-cov}
 \Sigma_\lambda(i,j)=
 \begin{cases}
 1,&i=j,\\
 \lambda/d,&\{i,j\}=\{0,\ell\},\\
 (\lambda^2-d)/[d(d-1)],&i,j\geq1,\ i\ne j.
 \end{cases}
\end{equation}
The star lies in $\{x:\sum_{i=1}^d x_i=\lambda x_0\}$.
The covariance identities are exact at separated eigenvalues too;
joint Gaussianity at separated limiting energies does not follow from them.

\section{A normalized local version for eigenvectors}\label{sec:normalized}
Let $G_N$ now be a uniform simple $d$-regular graph. The only randomness
in probabilities below is the graph; all local laws remain conditional
sampling laws. Define, with an empty supremum interpreted as zero,
\begin{equation}\label{eq:M}
 M_N(K)=\sup_{\substack{Au=\lambda u,\ \lambda\in K\\\|u\|_2=1}}
          N\sum_{v=1}^N u(v)^4.
\end{equation}
The assertion $M_N(K)=O_{\PP}(1)$ means: for every $\alpha>0$, there is
$C_\alpha<\infty$ with
$\limsup_N\PP(M_N(K)>C_\alpha)<\alpha$.

\begin{theorem}[Uniform normalized Gaussian neighborhoods]\label{thm:main}
Assume $M_N(K)=O_{\PP}(1)$. Fix $k\geq1$, finite
$F\subset V(T_d)$, and a deterministic sequence $\delta_N\downarrow0$.
Then
\begin{equation}\label{eq:main}
 \sup_{\substack{(u_1,\ldots,u_k)\ \mathrm{orthonormal}\ \mathrm{eigenvectors}\\
 \lambda_a\in K,\ \max_{a,b}|\lambda_a-\lambda_b|\leq\delta_N}}
 W_2\bigl(\Law_\phi(X^F),\gamma^F_{\boldsymbol\lambda}\bigr)
 \xrightarrow{\PP}0.
\end{equation}
Here $W_2$ uses the Euclidean cost in all $k|F|$ physical coordinates.
All real orthonormal choices inside multiple eigenspaces are included.
No rate is asserted. For $k=1$ the energy-diameter condition is vacuous.
\end{theorem}

\begin{lemma}[Fourth moments prevent second-moment loss]\label{lem:UI}
For $m=k|F|$, on $\{M_N(K)\leq C\}$,
\begin{equation}\label{eq:UI}
 \E_\phi\|X^F\|^4\leq Cm^2,\qquad
 \E_\phi[\|X^F\|^2\one_{\{\|X^F\|>L\}}]\leq Cm^2/L^2.
\end{equation}
Consequently every coordinate square and every coordinate product is
uniformly integrable along any such sequence of graphs and tuples.
\end{lemma}
\begin{proof}
Uniformity of $\phi(x)$ gives
$\E_\phi|X_{a,x}|^4=N\sum_v u_a(v)^4\leq C$.
Now $(\sum_{\ell=1}^m z_\ell^2)^2\leq m\sum_{\ell=1}^m z_\ell^4$.
The tail estimate follows from
$\|z\|^2\one_{\{\|z\|>L\}}\leq L^{-2}\|z\|^4$.
Products are bounded in absolute value by one half of the sum of
the corresponding squares.
\end{proof}

\begin{proof}[Proof of Theorem~\ref{thm:main}]
\emph{Step 1: the imported Gaussian classification.}
Use \BS's Theorem 2.2 and its explicit multichannel extension in Remark 1.
For a fixed tuple size these are uniform finite-graph conclusions, on
events whose probability tends to one. Along sequences of their good
graphs, with errors and eigenvalue diameters tending to zero, each weak
subsequential neighborhood limit is a centered jointly Gaussian
eigenvector process at the common limiting energy, possibly with lost
covariance. One may first take a subsequence on larger and larger
balls and then restrict to $F$.

The required common-energy condition is transparent for exact
orthonormal eigenvectors: for $\sum_a c_a^2=1$,
\begin{equation}\label{eq:combination}
 \left\|(A-\lambda_1)\sum_a c_a u_a\right\|_2^2
 =\sum_a c_a^2(\lambda_a-\lambda_1)^2\leq\delta_N^2.
\end{equation}
This identity explains applicability; the local multichannel
classification is imported from Remark 1. Scalar Cram\'er--Wold using
only such channel combinations would by itself test only coefficients
of the form $c_a t_x$ across channels and sites, rather than arbitrary
coefficients $t_{a,x}$. It is not a complete substitute for that remark.

\emph{Step 2: identify the covariance.}
Consider any sequence of the preceding good graphs on which
$M_N(K)\leq C$. After subselection, compactness of $K$ and the diameter
condition give $\lambda_a\to\lambda\in K$ for all $a$.
The fixed-dimensional laws are tight by their bounded second moments.
Lemma~\ref{lem:UI} lets their means and covariance entries pass to
each weak limit. Lemma~\ref{lem:cov} identifies the limiting covariance as
\[
 \E X_{a,x}X_{b,y}
 =\one_{a=b}\varphi_\lambda(\operatorname{dist}(x,y)).
\]
The centered joint Gaussian law is therefore exactly
$\gamma^F_{(\lambda,\ldots,\lambda)}$. Its channels are independent:
the characteristic function of a Gaussian with block diagonal
covariance factors over the blocks, including singular blocks.

\emph{Step 3: upgrade weak convergence to $W_2$.}
We recall the finite-dimensional implication with its relevant
justification. A Skorokhod coupling gives $Y_n\to Y$ almost surely.
Uniform integrability of $\|Y_n\|^2$, together with integrability of
$\|Y\|^2$, implies uniform integrability of
$\|Y_n-Y\|^2\leq2\|Y_n\|^2+2\|Y\|^2$. Hence
$\E\|Y_n-Y\|^2\to0$, which bounds $W_2^2$ from above.
Lemma~\ref{lem:UI} supplies that uniform integrability here.

The targets also converge in $W_2$. The entries of $\Ga_F(\lambda)$
are continuous in $\lambda$ by \eqref{eq:recurrence}. Couple the Gaussian
vectors using a single standard normal vector and the positive
semidefinite square roots of their covariance matrices. Continuity
of those square roots gives convergence in mean squared norm.
This argument handles changing singular supports without whitening.
The triangle inequality now gives the desired discrepancy tending to
zero along every such graph-and-tuple sequence.

\emph{Step 4: make the conclusion uniform and in probability.}
If \eqref{eq:main} failed, some subsequence would have discrepancy
larger than a fixed $\varepsilon>0$ with graph probability at least
$p>0$. Tightness of $M_N(K)$ gives $C$ with
$\limsup\PP(M_N(K)>C)<p/2$. A diagonal choice of the uniform
\BS\ events makes their accuracy tend to zero on a further
subsequence. For each large index choose a graph in their intersection
with $\{M_N(K)\leq C\}$ on which a bad tuple exists, and choose a tuple
with discrepancy at least $\varepsilon/2$. The intersection is nonempty
by the probability bounds. Steps 1--3 contradict this choice.
Thus the supremum tends to zero in probability. For a fixed $N$ the
graph sample space is finite, so measurability of this graph-valued
supremum causes no difficulty.
\end{proof}

\begin{remark}[What is actually needed]
The proof uses uniform integrability of root squares, transferred to
finitely many uniformly sampled sites. A bounded $2+\epsilon$ moment,
for any fixed $\epsilon>0$, can replace the fourth moment with the
corresponding tail estimate. A fourth moment tending to $3$, mixed
fourth-moment convergence, and a delocalization rate are unnecessary
for Theorem~\ref{thm:main}. A bound of the form
$\sqrt N\|u\|_\infty\leq(\log N)^C$ alone does not supply a uniform
tail modulus for the squares.
\end{remark}

\begin{corollary}[Almost eigenvectors with controlled tails]\label{cor:almost}
The same conclusion holds for any family of orthonormal zero-sum
tuples satisfying
\[
 \|(A-\lambda_a)u_a\|_2\leq\rho_N\to0,\qquad
 \lambda_a\in K,\qquad
 \max_{a,b}|\lambda_a-\lambda_b|\leq\delta_N\to0,
\]
provided their scalar empirical fourth moments are bounded uniformly
in probability over that family.
\end{corollary}
\begin{proof}
Write $r_a=(A-\lambda_a)u_a$. A unit channel combination has residual
at most $\sqrt k\rho_N+\delta_N$ at $\lambda_1$.
For each cover coordinate, the lifted error $\sqrt N r_a(\phi(x))$
has $L^2$ norm at most $\rho_N$. The covariance recurrence in
Lemma~\ref{lem:cov} therefore holds with errors bounded by $\rho_N$
before division by $d$ or $d-1$, by Cauchy--Schwarz. Induction to any
fixed distance gives an $O_{d,K,F}(\rho_N)$ covariance error.
Apply the joint almost-eigenvector version of the same \BS\ input
and repeat the proof. The tail condition is essential: the source
article constructs localized almost eigenvectors.
\end{proof}

\section{The precise interface to the bulk manuscript}\label{sec:interface}
There are two uses downstream. The proof at label
\texttt{rs:gaussian-limit} requires fixed tuples of edge marks at energies
tending to the same $E$. The variance and remote-tail arguments at
\texttt{mic:quadratic-covariance-bound} and \texttt{rs:tail} require
quantitative covariance control for quadratic marks, uniformly over
energies anywhere in a fixed bulk interval. The latter is an additional
input.

Here is a sufficient quantitative formulation smaller than the draft's
full mixed-fourth theorem. It involves at most two channels and only
edge coordinates.

\begin{assumption}[One-edge quadratic covariance]\label{ass:Q}
There are deterministic $a_N\downarrow0$, $C<\infty$, and graph events
$\Omega_N$ with $\PP(\Omega_N)\to1$ on which the following holds
simultaneously for all real unit bulk eigenvectors, and all orthonormal
pairs of them. On a uniform oriented edge $(v,w)$ set
\[
 \xi_i=(\sqrt N u_i(v),\sqrt N u_i(w)),\qquad
 B_i=B(\lambda_i)=
 \begin{pmatrix}1&\lambda_i/d\\\lambda_i/d&1\end{pmatrix},\qquad
 q_i^{rs}=\xi_{i,r}\xi_{i,s}-(B_i)_{rs}.
\]
For $r,s,t,u\in\{0,1\}$,
\begin{equation}\label{eq:Q}
 \left|\E_{\rm edge}q_i^{rs}q_j^{tu}\right|
 \leq C\one_{i=j}+a_N.
\end{equation}
The diagonal case means using the same vector twice; the off-diagonal
case means two orthonormal vectors, also when their eigenvalues coincide.
\end{assumption}

\begin{proposition}[Consequences of the reduced input]\label{prop:Q}
Assumption~\ref{ass:Q} implies Theorem~\ref{thm:main} and its
graph-measurable fixed-energy marked consequence below. It also gives,
for two independently sampled oriented edges and any finite set of
orthonormal bulk eigenvectors, the bound
\begin{equation}\label{eq:weighted}
 \E_*\left\|\sum_i h_i\bigl(y_i y_i^T-C_i\bigr)\right\|_{\rm F}^2
 \leq C'\sum_i|h_i|^2+C'a_N\left(\sum_i|h_i|\right)^2,
\end{equation}
where $y_i=(\xi_i^{(1)},\xi_i^{(2)})\in\R^4$,
$C_i=\operatorname{diag}(B_i,B_i)$, and $h_i$ are arbitrary complex
coefficients fixed with the graph. Constants depend only on $d,K,C$.
\end{proposition}
\begin{proof}
Since $\E_{\rm edge}\xi_{i,0}^2=1$, the diagonal case of
\eqref{eq:Q} gives
\[
 N\sum_v u_i(v)^4
 =\E_{\rm edge}\xi_{i,0}^4
 =1+\E_{\rm edge}(q_i^{00})^2\leq1+C+a_N.
\]
This holds simultaneously for all allowed vectors, so
$M_N(K)=O_{\PP}(1)$ and Theorem~\ref{thm:main} applies.

For the matrix estimate, let $D_i=y_i y_i^T-C_i$. Every entry is centered.
If the two quadratic factors $D_i^{pq},D_j^{rs}$ lie within the same
sampled edge, \eqref{eq:Q} applies. If each lies in a different edge,
independence of the edge samples and centering give zero. If three
of the four linear factors lie in one edge and one in the other,
the singleton has mean zero. In the remaining case each edge carries
one factor from channel $i$ and one from channel $j$. Their expectations
are exact cross-channel second moments, zero for $i\ne j$ by
Lemma~\ref{lem:cov} and bounded for $i=j$. Thus
\[
 |\E_*D_i^{pq}D_j^{rs}|\leq C'\one_{i=j}+a_N
\]
after adjusting constants. Expand the squared absolute value of each
entry of $\sum_i h_iD_i$, bound with $|h_i||h_j|$, and sum over the
sixteen entries to obtain \eqref{eq:weighted}.
\end{proof}

\paragraph{How the draft supplies this hypothesis.}
The draft's moment theorem at \texttt{th:stars} (also labeled
\texttt{en:fourth}) compares
every star quartic for up to four channels with its independent-Gaussian
Wick value at polynomial error $N^{-\chi}$. Restrict to one edge and
the channel pattern $i,i,j,j$. Exact second moments give
\[
 \E q_i^{rs}q_j^{tu}
 =\E\xi_{i,r}\xi_{i,s}\xi_{j,t}\xi_{j,u}-(B_i)_{rs}(B_j)_{tu}.
\]
For $i\ne j$ the Gaussian comparison is zero; for $i=j$ it is uniformly
bounded. Hence that theorem implies Assumption~\ref{ass:Q} with
$a_N=O(N^{-\chi})$. The full four-channel statement and exact diagonal
Wick approximation are stronger than this downstream use requires.
This is a reduction of the required conclusion, not an assertion that
the quantitative proof automatically becomes shorter.

\paragraph{Why a rate is needed.}
With $x_i=N\sigma(\lambda_i-E)$, $\sigma=\pi\rho_d(E)$, take
$h_i=\sigma/(x_i-w)$ over remote bulk marks. The second term of
\eqref{eq:weighted} is controlled by
\begin{equation}\label{eq:rate}
 C_Ba_N\left(\sum_{\lambda_i\in K}|x_i-w|^{-1}\right)^2.
\end{equation}
The draft's bound at \texttt{rs:absolute-sum} is $O_B(N^\varepsilon)$
for the sum, so this error vanishes if
$a_NN^{2\varepsilon}\to0$. A polynomial $a_N=N^{-\chi}$ suffices
when $2\varepsilon<\chi$. Qualitative $a_N=o(1)$ alone does not.
Independently proving a polylogarithmic bound on that sum would permit
a correspondingly weaker rate.

For $|x_i|>R\geq2\sup_{w\in B}|w|+1$, the diagonal term is controlled by
\[
 \sum_{|x_i|>R}|x_i-w|^{-2}
 \leq\frac{C_B}{R}\operatorname{Im}s_N(\ii R),\qquad
 s_N(w)=\sum_i(x_i-w)^{-1}.
\]
The draft's capped height-moment bounds then give the required $O_B(R^{-1})$
tail in expectation. Both the rate in \eqref{eq:rate} and those separate
spectral estimates are used; a fixed-tuple Gaussian limit cannot replace
them.

\begin{corollary}[Conditional Gaussian edge marks]\label{cor:marks}
Under Theorem~\ref{thm:main}'s hypothesis, fix a graph-measurable
orthonormal tuple of size $k$ with $\lambda_{a,N}\in K$ with probability
tending to one and $\lambda_{a,N}\to E\in K$ in probability.
Sample two oriented edges independently conditional on the
graph, and let $Y_N=(y_{1,N},\ldots,y_{k,N})$. Put
$C(E)=\operatorname{diag}(B(E),B(E))$ and let
$G_1,\ldots,G_k$ be independent $N(0,C(E))$ vectors. For every fixed
bounded Lipschitz $f$,
\begin{equation}\label{eq:conditional}
 \E_*f(Y_N)-\E f(G_1,\ldots,G_k)\longrightarrow0
 \quad\text{in probability and in }L^1.
\end{equation}
For any graph-measurable $H_N$ with $|H_N|\leq1$,
\begin{equation}\label{eq:stable}
 \E[H_N f(Y_N)]-\E[H_N]\E f(G_1,\ldots,G_k)\longrightarrow0.
\end{equation}
Consequently, whenever a graph-measurable state $S_N$ has a joint
subsequential limit with these marks, the marks in that limit have the
product Gaussian law and are independent of the limiting state.
\end{corollary}
\begin{proof}
Convergence of finitely many energies to $E$ permits a deterministic
$\delta_N\downarrow0$ dominating their diameter with probability tending
to one. Projection of Theorem~\ref{thm:main} to an oriented edge gives
conditional convergence for one edge. Two conditionally independent
samples give the product law and hence \eqref{eq:conditional}, with
covariances tending to $C(E)$. Bounded convergence in probability gives
$L^1$ convergence. Multiply the conditional error by $H_N$ to obtain
\eqref{eq:stable}. Bounded continuous product tests identify independence
in any joint weak limit.
\end{proof}
This formulation permits the graph-dependent spectral selection and
padding of the tuple used at \texttt{rs:gaussian-limit}, including
multiplicities and arbitrary measurable sign choices. The tuple size
is fixed before $N\to\infty$. Conditional independence of two edge
samples is exact even though their endpoint sets can overlap.

\section{What a quantitative extension of the article must add}\label{sec:quantitative}
For normalized local Gaussianity the preceding proof is sufficient
once scalar tail control is available. To establish a polynomial
version of Assumption~\ref{ass:Q}, one needs new quantitative work.
The supplied bulk draft proposes the following route, which refines
the original entropy argument rather than following from its stated
theorems.

\paragraph{1. Counting with correlated Gaussian noise.}
On a common counting event, let $x$ be a bounded graph signal and
$\zeta\sim N(0,Q)$ a nondegenerate, possibly graph-dependent, Gaussian
field. Control the difference between realized and averaged edge-color
entropies using Gaussian Poincar\'e and quantization. The proposed
channel bound is
\[
 h((x+\zeta)_S)-\frac d2h((x+\zeta)_e)
 \geq\frac1Nh(\zeta)-\text{polynomial error}.
\]
Uniformity over graph-dependent $Q$ is justified by a counting event
uniform over all colorings; it is a substantive part of this extension.

\paragraph{2. Resolvent calibration.}
Use independent channel fields with
\[
 Q_a=\sigma_{\lambda_a}^{-1}
       \operatorname{Im}(A-\lambda_a-\ii\eta)^{-1},\qquad
 \sigma_\lambda=\pi\rho_d(\lambda),\quad\eta=N^{-c}.
\]
The local tangent covariance approaches the Gaussian-star covariance.
For unit normal $(-\lambda,1,\ldots,1)/\sqrt{d+\lambda^2}$ its variance
is $\eta/[\sigma_\lambda(d+\lambda^2)]$ to small relative error.
The singular normal entropy cancels the corresponding term in the
global noise entropy. With
$L_d(\lambda)=\int\rho_d(t)\log|t-\lambda|\dd t$, the remaining identity is
\[
 \left[\frac12\log\frac{2\pi\ee\eta}{\sigma_\lambda}-L_d(\lambda)\right]
 -\frac12\log\frac{2\pi\ee\eta}{\sigma_\lambda(d+\lambda^2)}
 =\frac12\log(d+\lambda^2)-L_d(\lambda).
\]
This produces the Gaussian-star entropy benchmark without first
assuming local Gaussianity of the eigenvector signal.

\paragraph{3. High-noise expansion and fourth chaos.}
After whitening, let $Y$ be centered, isotropic, symmetry-admissible,
and $\|Y\|\leq B$, where $B\geq2$. Write $f_s$ for the density relative to standard
Gaussian measure of $\sqrt sY+\sqrt{1-s}G$, and
$f_s-1=\sum_{j\geq3}s^{j/2}g_j$. The entropy deficit has the proposed
expansion, valid for $0<s\leq cB^{-2}$,
\[
 \Delta_s=\frac12\mathcal Q(f_s-1)+O(B^3s^{9/2}),\qquad
 \mathcal Q(h)=\|h\|_2^2-\frac d2\|\mathcal C_e h\|_2^2.
\]
Here $\mathcal C_e$ is Gaussian conditional expectation onto one edge.
Nonnegativity on admissible Hermite chaoses and strict fourth-chaos
coercivity give
$\mathcal Q(f_s-1)\geq\kappa s^4\|g_4\|_2^2$.
These are analytic and finite-dimensional geometric lemmas to verify
separately, including their multichannel versions at arbitrary bulk
energies. Scalar near-energy Gaussianity is not a proof of them.

\paragraph{4. Extract the rate.}
If monotonicity gives $0\leq\Delta_s\leq\epsilon$ for
$0<s\leq1/2$, the preceding estimates yield
\[
 \|g_4\|_2^2\leq C(\epsilon s^{-4}+B^3s^{1/2}).
\]
For $\epsilon\leq cB^{-6}$, choose
$s=(\epsilon/B^3)^{2/9}$ to obtain
\[
 \|g_4\|_2\leq CB^{4/3}\epsilon^{1/18}.
\]
The difference of expectations of a fixed homogeneous quartic is its
inner product with $g_4$, since degrees zero and two already match.
Thus $B=(\log N)^C$ and $\epsilon=N^{-\gamma}$ give any smaller
polynomial exponent $\chi<\gamma/18$. Restricting the resulting
quartics to two channels on one edge proves the reduced input.

The displayed cancellation and balancing exponents are consistent.
This note does not provide a complete independent verification of the
draft's counting, local-law, entropy-expansion, and coercivity proofs.
Their statements remain explicitly separate quantitative dependencies
of the bulk application.

\section{Limits of the conclusion and a Lean proof structure}\label{sec:lean}
Two simple measure examples explain why the hypotheses cannot be
silently removed. They are logical counterexamples about moment
convergence, not constructions of graph eigenvectors.

For $0\leq s<1$, let
\[
 \mu_n=(1-n^{-2})N(0,s^2)
       +\frac{n^{-2}}2(\delta_{a_n}+\delta_{-a_n}),\qquad
 a_n^2=n^2(1-s^2)+s^2.
\]
These laws have mean zero and variance exactly one, but converge weakly
to $N(0,s^2)$. Second moments escape in a rare tail. This is the
variance-loss issue addressed by Lemma~\ref{lem:UI}.

Even a bounded fourth moment does not ensure fourth-moment convergence.
For $n\geq2$ set
\[
 \nu_n=(1-n^{-4})N(0,s_n^2)
       +\frac{n^{-4}}2(\delta_n+\delta_{-n}),\qquad
 s_n^2=\frac{1-n^{-2}}{1-n^{-4}}.
\]
Then $\int x^2\dd\nu_n=1$, $\nu_n\to N(0,1)$ in $W_2$, and
\[
 \int x^4\dd\nu_n
 =3\frac{(1-n^{-2})^2}{1-n^{-4}}+1\longrightarrow4.
\]
Thus the normalization proof cannot establish the draft's mixed
fourth-moment theorem. Also, Gaussian marginals with zero covariance
need not be independent: $(Z,SZ)$, for independent standard Gaussian
$Z$ and Rademacher $S$, has equal absolute coordinates. We use
\emph{joint} Gaussianity to deduce independence.

\paragraph{Recommended formalization boundary.}
The first Lean target should be the finite-dimensional normalization
and interface argument, with the published qualitative theorem supplied
as an explicit hypothesis. The following order separates finite
algebra from the deep entropy classification.
\begin{center}
\small
\begin{tabular}{>{\raggedright\arraybackslash}p{.28\linewidth}>{\raggedright\arraybackslash}p{.46\linewidth}>{\raggedright\arraybackslash}p{.14\linewidth}}
\toprule
Suggested lemma/module & Mathematical contract & Status here\\
\midrule
\texttt{CoverSampling} & A finite restriction of the uniform covering
law; uniform vertex marginals and edge reversal. & Defined / used\\[3pt]
\texttt{EigenCovariance} & Zero means, Gram identity, and covariance
recurrence \eqref{eq:cov}; finite star formulas first. & Proved\\[3pt]
\texttt{MomentTail} & Root fourth bound implies vector fourth and
square-tail bounds \eqref{eq:UI}. & Proved\\[3pt]
\texttt{BSJointLimit} & Uniform joint weak Gaussian classification on
fixed neighborhoods at a common limiting energy. & Imported \cite{BS}\\[3pt]
\texttt{NormalizeLimit} & Weak joint Gaussian limit plus square-tail
control and covariance identities fixes its normalized product law.
& Proved\\[3pt]
\texttt{WassersteinUpgrade} & Square uniform integrability upgrades
weak convergence; PSD-root coupling gives target continuity. & Proved\\[3pt]
\texttt{EdgeQuadratics} & Assumption~\ref{ass:Q} implies scalar fourth
control, two-edge covariance, and \eqref{eq:weighted}. & Proved\\[3pt]
\texttt{ConditionalMarks} & Quenched bounded-Lipschitz convergence
implies \eqref{eq:stable} and independence from a graph state.
& Proved\\
\bottomrule
\end{tabular}
\end{center}

In a first formalization, parameterize \texttt{NormalizeLimit} by a
named proposition asserting \texttt{BSJointLimit}; do not disguise this
dependency as a completed formal proof. Likewise, expose
Assumption~\ref{ass:Q} and its rate as inputs to
\texttt{EdgeQuadratics}. No Lean files or machine-checked results are
claimed by this note. A second stage could formalize the entropy and
fourth-chaos route in Section~\ref{sec:quantitative}.

For the manuscript, the clean dependency is therefore:
one-edge quantitative quadratic covariance has two consequences:
\begin{enumerate}
\item It gives root square-tail control. Combined with the \BS\ joint
local theorem, this yields normalized Gaussian marks.
\item It gives the weighted quadratic variance bound. Combined with the
separate spectral estimates, this yields remote marked-tail control.
\end{enumerate}
This is the eigenvector input package. The switching identities,
microscopic tightness, loop hierarchy, and point-process identification
remain the separate downstream arguments of the bulk manuscript.

\begin{thebibliography}{9}
\bibitem{BS}
\'A. Backhausz and B. Szegedy,
\emph{On the almost eigenvectors of random regular graphs},
Ann. Probab.\ \textbf{47} (2019), no.\ 3, 1677--1725.
\href{https://doi.org/10.1214/18-AOP1294}{doi:10.1214/18-AOP1294};
\href{https://arxiv.org/abs/1607.04785}{arXiv:1607.04785}.
The supplied published PDF is the source for theorem numbering.
\bibitem{Bulk}
Y. He, J. Huang, and X. Wang,
\emph{Bulk universality of random regular graphs}, supplied TeX draft,
\texttt{bulk-universality-point-process.tex}, file dated November 3, 2026.
Relevant labels: \texttt{th:stars}, \texttt{en:qualitative},
\texttt{ch:channel}, \texttt{en:deficit}, \texttt{en:smoothing},
\texttt{en:quartic-coercivity}, \texttt{mic:quadratic-covariance-bound},
\texttt{rs:absolute-sum}, \texttt{rs:tail}, \texttt{rs:gaussian-limit}.
\end{thebibliography}
\end{document}
